% GENERATED FILE. Edit canonical lessons, then run npm run generate:books.
% canonical-source-hash: fnv1a64:764fa6103c194090
% canonical-lessons: IT-C322-fuochi-d-artificio, IT-C322-botto, IT-C322-falo, IT-C322-lanterna, IT-C322-danzare

\chapter{Fuochi e falò}
\label{ch:it-fuochi-e-falo}
\hlchaptermodality{322}

\begin{chapteropening}
\textbf{By the end of this chapter:} \emph{I can describe a night of fireworks, bangs, bonfires and lanterns, with dancing.}

Complete the last lesson of chapter 322: I can describe a night of fireworks, bangs, bonfires and lanterns, with dancing.
\end{chapteropening}

\section[i fuochi d'artificio]{i fuochi d'artificio — the fireworks}

\label{lesson:IT-C322-fuochi-d-artificio}



\begin{quote}
Before the new one: say the Italian for to decorate, then the Italian for the tombola.
\end{quote}

\subsection*{You'll want to know: i fuochi d'artificio}
\textbf{i fuochi d'artificio} — \textquotedblleft{}the fireworks\textquotedblright{}.

\begin{quote}
A mezzanotte ci sono i fuochi d'artificio.
\end{quote}

\begin{quote}
\emph{At midnight there are fireworks.}
\end{quote}

\emph{Artificio} — a skilful device.

\begin{cousinweb}
\emph{fuochi} is the plural of \emph{fuoco} \textquotedblleft{}fire\textquotedblright{}, from Latin \emph{focus} \textquotedblleft{}hearth\textquotedblright{} (Lewis \& Short); \emph{artificio} comes from Latin \emph{artificium} \textquotedblleft{}craft, skill\textquotedblright{}.
\end{cousinweb}

\begin{grammarlens}[title={What you've built}]
One more word for telling, explaining and describing at B1.
\end{grammarlens}

\subsection*{Your turn}
\begin{itemize}
\raggedright
  \item \emph{Say it:} \emph{i fuochi d'artificio}
  \item \emph{Say it:} \emph{i fuochi d'artificio}, once more
  \item \emph{Say it:} a sentence of your own with \emph{i fuochi d'artificio}
  \item \emph{Recall:} say the Italian for to decorate, then the Italian for the tombola, then say \emph{i fuochi d'artificio} again
\end{itemize}

\begin{tcolorbox}[breakable,colback=teal!4,colframe=teal!35!black,title={Before you move on}]
What does i fuochi d'artificio mean? (\textquotedblleft{}The fireworks\textquotedblright{}.) Say it once more.
\end{tcolorbox}
\section[il botto]{il botto — the firework bang; the crash}

\label{lesson:IT-C322-botto}



\begin{quote}
Before the new one: say the Italian for the tombola, then the Italian for the fireworks.
\end{quote}

\subsection*{You'll want to know: il botto}
\textbf{il botto} — \textquotedblleft{}the firework bang; the crash\textquotedblright{}.

\begin{quote}
A fine anno si sentono botti dappertutto.
\end{quote}

\begin{quote}
\emph{At New Year you hear bangs everywhere.}
\end{quote}

\emph{Di botto} — suddenly. \emph{I botti} are firecrackers.

\begin{cousinweb}
\emph{botto} is usually explained as a sound word that imitates a bang (EtymDB has no entry).
\end{cousinweb}

\begin{grammarlens}[title={What you've built}]
One more word for telling, explaining and describing at B1.
\end{grammarlens}

\subsection*{Your turn}
\begin{itemize}
\raggedright
  \item \emph{Say it:} \emph{il botto}
  \item \emph{Say it:} \emph{il botto}, once more
  \item \emph{Say it:} a sentence of your own with \emph{il botto}
  \item \emph{Recall:} say the Italian for the tombola, then the Italian for the fireworks, then say \emph{il botto} again
\end{itemize}

\begin{tcolorbox}[breakable,colback=teal!4,colframe=teal!35!black,title={Before you move on}]
What does il botto mean? (\textquotedblleft{}The firework bang; the crash\textquotedblright{}.) Say it once more.
\end{tcolorbox}
\section[il falò]{il falò — the bonfire}

\label{lesson:IT-C322-falo}



\begin{quote}
Before the new one: say the Italian for the fireworks, then the Italian for the firework bang.
\end{quote}

\subsection*{You'll want to know: il falò}
\textbf{il falò} — \textquotedblleft{}the bonfire\textquotedblright{}.

\begin{quote}
Sulla spiaggia abbiamo acceso un falò.
\end{quote}

\begin{quote}
\emph{We lit a bonfire on the beach.}
\end{quote}

Stress on the last syllable; plural unchanged.

\begin{cousinweb}
\emph{falò} comes from Greek \emph{phanós} \textquotedblleft{}a lantern, a torch\textquotedblright{} (EtymDB), probably through a Byzantine Greek form.
\end{cousinweb}

\begin{grammarlens}[title={What you've built}]
One more word for telling, explaining and describing at B1.
\end{grammarlens}

\subsection*{Your turn}
\begin{itemize}
\raggedright
  \item \emph{Say it:} \emph{il falò}
  \item \emph{Say it:} \emph{il falò}, once more
  \item \emph{Say it:} a sentence of your own with \emph{il falò}
  \item \emph{Recall:} say the Italian for the fireworks, then the Italian for the firework bang, then say \emph{il falò} again
\end{itemize}

\begin{tcolorbox}[breakable,colback=teal!4,colframe=teal!35!black,title={Before you move on}]
What does il falò mean? (\textquotedblleft{}The bonfire\textquotedblright{}.) Say it once more.
\end{tcolorbox}
\section[la lanterna]{la lanterna — the lantern}

\label{lesson:IT-C322-lanterna}



\begin{quote}
Before the new one: say the Italian for the firework bang, then the Italian for the bonfire.
\end{quote}

\subsection*{You'll want to know: la lanterna}
\textbf{la lanterna} — \textquotedblleft{}the lantern\textquotedblright{}.

\begin{quote}
Le lanterne illuminano la piazza.
\end{quote}

\begin{quote}
\emph{The lanterns light up the square.}
\end{quote}

In some ports, the lighthouse is called \emph{la lanterna}.

\begin{cousinweb}
\emph{lanterna} is inherited from Latin \emph{lanterna} (EtymDB), from Greek \emph{lampt\'{\={e}}r} \textquotedblleft{}a torch\textquotedblright{}.
\end{cousinweb}

\begin{grammarlens}[title={What you've built}]
One more word for telling, explaining and describing at B1.
\end{grammarlens}

\subsection*{Your turn}
\begin{itemize}
\raggedright
  \item \emph{Say it:} \emph{la lanterna}
  \item \emph{Say it:} \emph{la lanterna}, once more
  \item \emph{Say it:} a sentence of your own with \emph{la lanterna}
  \item \emph{Recall:} say the Italian for the firework bang, then the Italian for the bonfire, then say \emph{la lanterna} again
\end{itemize}

\begin{tcolorbox}[breakable,colback=teal!4,colframe=teal!35!black,title={Before you move on}]
What does la lanterna mean? (\textquotedblleft{}The lantern\textquotedblright{}.) Say it once more.
\end{tcolorbox}
\section[danzare]{danzare — to dance (formal)}

\label{lesson:IT-C322-danzare}



\begin{quote}
Before the new one: say the Italian for the bonfire, then the Italian for the lantern.
\end{quote}

\subsection*{You'll want to know: danzare}
\textbf{danzare} — \textquotedblleft{}to dance (formal)\textquotedblright{}.

\begin{quote}
Abbiamo danzato intorno al fuoco.
\end{quote}

\begin{quote}
\emph{We danced around the fire.}
\end{quote}

A literary word for \emph{ballare}.

\begin{cousinweb}
\emph{danzare} was borrowed from French \emph{danser} (EtymDB), whose deeper origin is not settled.
\end{cousinweb}

\begin{grammarlens}[title={What you've built}]
One more word for telling, explaining and describing at B1.
\end{grammarlens}

\subsection*{Your turn}
\begin{itemize}
\raggedright
  \item \emph{Say it:} \emph{danzare}
  \item \emph{Say it:} \emph{danzare}, once more
  \item \emph{Say it:} a sentence of your own with \emph{danzare}
  \item \emph{Recall:} say the Italian for the bonfire, then the Italian for the lantern, then say \emph{danzare} again
\end{itemize}

\begin{tcolorbox}[breakable,colback=teal!4,colframe=teal!35!black,title={Before you move on}]
What does danzare mean? (\textquotedblleft{}To dance (formal)\textquotedblright{}.) Say it once more.
\end{tcolorbox}
